\appendix
\section{Proofs of Theoretical Results}
\label{app:proofs}
Before we prove Theorem~\ref{thm:detectability} in Appendix~\ref{proof_theorem}, we need the following intermediate lemmas on mean-field and concentration lemma and the spectral norm of the difference between adjacency matrix and its mean.

\subsection{Mean-field and Concentration Lemmas}

\begin{lemma}[Mean-field $2\times 2$ reduction]
\label{lem:meanfield}

Let $M=\begin{pmatrix} q_{II} & q_{IO}\\ q_{OI} & q_{OO}\end{pmatrix}$ denote the mean-field block transition matrix.
Then, for a seed node $s_0\in S^\star$,
\[
\mathbb{E}[r(S^\star)] = (1-\alpha)\sum_{k=0}^\infty \alpha^k (M^k)_{11}
= (1-\alpha)\big[(I-\alpha M)^{-1}\big]_{11}.
\]
\end{lemma}

\begin{proof}
The PPR vector satisfies
\[
\mathbf{r} = (1-\alpha)\sum_{k\ge0}\alpha^k P^k \mathbf{e}_{s_0}.
\]
The total PPR mass on the planted block $S^\star$ is
\[
r(S^\star) = \mathbf{1}_{S^\star}^\top \mathbf{r}
= (1-\alpha)\sum_{k\ge0}\alpha^k 
\mathbf{1}_{S^\star}^\top P^k \mathbf{e}_{s_0}.
\]
Taking expectations over the SBM ensemble and exploiting node exchangeability within each block,
the expected probability that a random walk is in block $I$ after $k$ steps is governed by the 2-state Markov chain with transition matrix $M$. Hence,
\[
\mathbb{E}\!\left[\mathbf{1}_{S^\star}^\top P^k \mathbf{e}_{s_0}\right]
= (M^k)_{11}.
\]
Summing the geometric series yields
\[
\mathbb{E}[r(S^\star)] = (1-\alpha)\sum_{k\ge0}\alpha^k (M^k)_{11}
= (1-\alpha)\big[(I-\alpha M)^{-1}\big]_{11},
\]
which completes the proof.
\end{proof}



Following Lemma 10 in~\parencite{avrachenkov2015pagerank}, we provide lemma~\ref{lem:concentration} below as the supplement to the assumption in Theorem~\ref{thm:detectability}.
\begin{lemma}[Degree and adjacency concentration]
\label{lem:concentration}
Let $d_v$ denote the degree of node $v$ in $\mathcal{G}$. There exist constants $C_1,C_2>0$ such that, with probability at least $1-n^{-c}$:
\begin{enumerate}
\item[(i)] (\textbf{Degrees}) For every node $v$,
\[
    |d_v-\mathbb{E}[d_v]| \le C_1\sqrt{\mathbb{E}[d_v]\log n}.
\]
\item[(ii)] (\textbf{Adjacency spectral norm})
\[
    \|A-\bar A\| \le C_2\sqrt{n p_{\max}\log n}, \quad p_{\max} = \max\{p_{\mathrm{in}},p_{\mathrm{out}}\}.
\]
\end{enumerate}
\end{lemma}

\begin{proof}
We divide the proof into two parts: (i) concentration of node degrees and (ii) concentration of the adjacency matrix in spectral norm.

\noindent(i) Given a node $v$, its degree can be written as $d_v = \sum_{u\neq v} A_{vu}$, where $\{A_{vu}\}_{u\neq v}$ are independent Bernoulli random variables with mean either $p_{\mathrm{in}}$ or $p_{\mathrm{out}}$ depending on the community memberships. Let $\mu_v = \mathbb{E}[d_v]$ denote its expected degree. Each $A_{vu}\in[0,1]$, and hence $d_v$ is a sum of $(n-1)$ independent bounded random variables.
By Bernstein’s inequality~\parencite{govil1987bernstein}, for any $t>0$, we have
\[
    \Pr(|d_v - \mu_v| > t)
    \le 2\exp\!\left(-\frac{t^2}{2(\mu_v+t/3)}\right).
\]
Setting $t=C_1\sqrt{\mathbb{E}[ \deg(v) ]\log n}$, we have 
\[
    \Pr(|d_v - \mu_v| > t) \le 2\exp\!\left(-\frac{C_1^2 \mu_v \log n}{2(\mu_v + \tfrac{1}{3}C_1\sqrt{\mu_v\log n})} \right) =
    2\exp\!\left( -\frac{C_1^2\log n}{2(1 + \tfrac{C_1}{3}\sqrt{\tfrac{\log n}{\mu_v}})} \right).
\]
When $\mu_v \gtrsim \log n$ (as assumed in standard dense or semi-sparse regimes), the denominator in the exponent is $O(1)$, so for sufficiently large $C_1$ we have
\[
    \Pr(|d_v - \mu_v| > t) \le 2n^{-c-1}.
\]
\noindent Applying a union bound over $n$ nodes gives
\[
    \Pr\! \left(\exists v:\, |d_v - \mu_v| > C_1\sqrt{\mu_v\log n}\right) \le 2n^{-c}.
\]
Hence, with probability at least $1-n^{-c}$,
\[
    |d_v-\mathbb E[d_v]| \le C_1\sqrt{\mathbb E[d_v]\log n} \quad\text{for all }v\in[n],
\]
which proves part (i).



\noindent(ii) For the spectral bound, $A-\bar A=\sum_{i<j} X_{ij}$, where $X_{ij}$ is the symmetric zero-mean random matrix with $(i,j)$ and $(j,i)$ entries equal to $A_{ij}-\mathbb{E}[A_{ij}]$. To apply the Matrix Bernstein inequality~\parencite{tropp2012user}, we compute the matrix variance parameter:
\[
    \sigma^2 := \left\|\sum_{i<j}\mathbb E[X_{ij}^2]\right\|.
\]
Observe that $X_{ij}^2$ is diagonal with entries 
$\mathbb E[(A_{ij}-\mathbb E[A_{ij}])^2] = \mathrm{Var}(A_{ij}) \le p_{\max}$ 
on rows $i$ and $j$. Therefore, for some absolute constant $C>0$, we have
\[
    \sigma^2 \le \left\| \sum_{j\neq i} p_{\max} e_i e_i^\top \right\| \le  C\,n p_{\max}.
\] 
The maximum single-summand norm satisfies $\|X_{ij}\|\le 1$.  The Matrix Bernstein inequality yields, for all $t>0$,
\begin{equation}
    \Pr(\|A-\bar A\| \ge t) \le  n \exp\!\left(-\frac{t^2/2}{\sigma^2 + t/3} \right).
    \label{eq:matrixbern}
\end{equation}
Setting $t = C_2\sqrt{n p_{\max}\log n} $ and using $\sigma^2 \le C n p_{\max}$ and $t\ll n p_{\max}$, the exponent in~\eqref{eq:matrixbern} satisfies
\[
\frac{t^2}{\sigma^2 + t/3}
\gtrsim
\frac{C_2^2 n p_{\max} \log n}{C n p_{\max} + C_2\sqrt{n p_{\max}\log n}/3}
\gtrsim c' C_2^2 \log n.
\]
Hence, for sufficiently large $C_2$, the probability in~\eqref{eq:matrixbern} is at most $n^{-c}$, proving that with probability at least $1-n^{-c}$,
\[
\|A-\bar A\| \le C_2\sqrt{n p_{\max}\log n}.
\]\noindent 
This completes the proof of Lemma~\ref{lem:concentration}.
\end{proof}




\begin{lemma}[Transition operator perturbation]
\label{lem:transition}
Under the high-probability event of Lemma~\ref{lem:concentration} and assuming $\min_i d_i \ge c_0 n p_{\mathrm{out}}$ for some constant $c_0>0$,
% and assuming $\min_i d_i \ge c_0 \log n$ for some constant $c_0 > 0$, 
there exists a constant $C>0$ such that
\begin{equation}
    \label{eq:epsinf}
    \|P-\bar P\| \le \varepsilon_P := C\frac{\sqrt{n p_{\max}\log n}}{n p_{\mathrm{out}}},
\end{equation}
where $\bar P$ is the mean-field transition matrix.
\end{lemma}


\begin{proof}
We first decompose $P - \bar P$:
\[
P - \bar P = D^{-1}A - \bar D^{-1}\bar A = D^{-1}(A-\bar A) + (D^{-1}-\bar D^{-1})\bar A.
\]
Hence, we have
\[
\|P-\bar P\| \le \|D^{-1}\|\|A-\bar A\| + \|D^{-1}-\bar D^{-1}\|\|\bar A\|.
\]

\noindent
\textbf{(1) Bounding $\|D^{-1}\|\|A-\bar A\|$:}  
By Lemma~\ref{lem:concentration}, $\min_i d_i \ge \frac{1}{2}c_0 n p_{\mathrm{out}}$ with high probability, so $\|D^{-1}\|\le \frac{2}{c_0 n p_{\mathrm{out}}}$.  Thus, we have
\[
\|D^{-1}\|\|A-\bar A\|
\le C_1'\frac{\sqrt{n p_{\max}\log n}}{n p_{\mathrm{out}}}.
\]

\noindent
\textbf{(2) Bounding $\|D^{-1}-\bar D^{-1}\|\|\bar A\|$:}  Using $D^{-1}-\bar D^{-1} = -\bar D^{-1}(D-\bar D)D^{-1}$, we obtain
\[
\|D^{-1}-\bar D^{-1}\|
\le \|\bar D^{-1}\|\|D-\bar D\|\|D^{-1}\|
\le C_2'\frac{\sqrt{n p_{\max}\log n}}{(n p_{\mathrm{out}})^2}.
\]
Since $\|\bar A\|\le C_3' n p_{\max}$, the second term is
\[
\|D^{-1}-\bar D^{-1}\|\|\bar A\|
\le C_4'\frac{\sqrt{n p_{\max}\log n}}{n p_{\mathrm{out}}}
\frac{p_{\max}}{p_{\mathrm{out}}}.
\]
Assuming $p_{\mathrm{in}},p_{\mathrm{out}}$ are of the same order,  the scaling matches the first term. 

\noindent Combining the results from (1) and (2), we have 
\begin{align}
    \notag \|P-\bar P\| &\le \|D^{-1}\|\|A-\bar A\| + \|D^{-1}-\bar D^{-1}\|\|\bar A\| \\
    \notag &\le (C_1' +C_4'\frac{p_{\max}}{p_{\mathrm{out}}}) \frac{\sqrt{n p_{\max}\log n}}{n p_{\mathrm{out}}}  \\
    \notag &= C\frac{\sqrt{n p_{\max}\log n}}{n p_{\mathrm{out}}} =  \varepsilon_P
\end{align}
which completes the proof.
\end{proof}






\begin{lemma}[Resolvent / Neumann propagation]
\label{lem:resolvent}
Assume $\alpha\|\bar P\|\le 1-\delta$ for some $\delta\in(0,1]$ and $\alpha\varepsilon_P \le \delta/2$. Then, with high probability,
\[
    \|(I-\alpha P)^{-1} - (I-\alpha \bar P)^{-1}\| \le \frac{2\alpha\varepsilon_P}{\delta^2}.
\]
\end{lemma}

\begin{proof}
By the resolvent identity,
\[
    (I-\alpha P)^{-1} - (I-\alpha \bar P)^{-1} = (I-\alpha P)^{-1}\alpha(P-\bar P)(I-\alpha \bar P)^{-1}.
\]
We take the spectral norm on both sides:
\[
    \|(I-\alpha P)^{-1} - (I-\alpha \bar P)^{-1}\| \le \|(I-\alpha P)^{-1}\| \cdot \alpha\|P-\bar P\| \cdot \|(I-\alpha \bar P)^{-1}\|.
\]
The assumption $\alpha\|\bar P\|\le 1-\delta$ implies that all eigenvalues of $\alpha\bar P$ lie within the unit disk of radius $(1-\delta)$. Therefore, the smallest singular value of $(I - \alpha\bar P)$ is at least $\delta$, and hence
\[
\|(I - \alpha \bar P)^{-1}\| \le \frac{1}{\delta}.
\]
By the triangle inequality and Lemma~\ref{lem:transition}, $\|P\| \le \|\bar P\| + \|P - \bar P\| \le \|\bar P\| + \varepsilon_P$.
Then
\[
\alpha\|P\| \le \alpha\|\bar P\| + \alpha\varepsilon_P 
\le (1 - \delta) + \frac{\delta}{2} = 1 - \frac{\delta}{2}.
\]
Hence, the spectral radius of $\alpha P$ is at most $(1 - \frac{\delta}{2})$, 
so $(I - \alpha P)$ is invertible, and by the Neumann series expansion
\[
(I - \alpha P)^{-1} = \sum_{k=0}^{\infty} (\alpha P)^k, \qquad \text{for } \|\alpha P\|<1.
\]
Taking norms gives
\[
\|(I - \alpha P)^{-1}\| \le \frac{1}{1 - \alpha\|P\|} 
\le \frac{1}{\delta/2} = \frac{2}{\delta}.
\]
Combining the above inequalities, we obtain
\[
    \|(I - \alpha P)^{-1} - (I - \alpha \bar P)^{-1}\| \le \frac{2}{\delta} \cdot \alpha\varepsilon_P \cdot \frac{1}{\delta} = \frac{2\alpha\varepsilon_P}{\delta^2}.
\]
which completes the proof. 
\end{proof}

% This lemma shows that the resolvent operator $(I - \alpha P)^{-1}$ is \emph{Lipschitz-stable} with respect to perturbations in $P$, with sensitivity scaling as $O(\alpha / \delta^2)$. Here $\delta$ represents the spectral gap controlling the distance of $\alpha\bar P$ from instability ($\|\alpha\bar P\| = 1$),  and $\varepsilon_P$ quantifies the deviation between $P$ and $\bar P$. When $\delta$ is small (i.e., the system is near the critical diffusion limit), the sensitivity of the Personalized PageRank operator grows as $1/\delta^2$.

\subsubsection{Proof of Theorem~\ref{thm:detectability}}
\begin{theorem}[Detectability of PPR under Planted Laundering Group]
Let $G\sim\mathrm{SBM}(n,s,p_{\mathrm{in}},p_{\mathrm{out}})$ with $p_{\mathrm{in}}>p_{\mathrm{out}}$ and assume that degrees concentrate around their expectations (i.e., Lemma~\ref{lem:concentration} holds). Let $\mathbf{r}$ be the Personalized PageRank vector seeded at a known malicious account $s_0\in S^\star$ with a constant $\alpha\in(0,1)$, and $\varepsilon_{\infty}$ denote the entrywise perturbation scale arising from stochastic fluctuations of $P$ around the mean-field transition matrix $\bar P$. There exist constants $C,c>0$ such that, with probability at least $1-n^{-c}$, if
\[
\Delta_{\mathrm{mean}} \;\ge\; C\cdot\varepsilon_{\infty},
\]
then
\begin{enumerate}
    \item the average PPR score on $S^\star$ exceeds that on $V\setminus S^\star$ by at least $\tfrac{1}{2}\Delta_{\mathrm{mean}}$;
    \item ordering nodes by normalized PPR $r(v)/d(v)$ and selecting the prefix with smallest conductance recovers a subset $\widehat S$ with $|\widehat S\cap S^\star|\ge\gamma s$ for some constant $\gamma\in(0,1)$.
\end{enumerate}
In the standard regime $p_{\mathrm{out}}\gtrsim \log n/n$,
the detectability condition simplifies to
\[
s\,(p_{\mathrm{in}}-p_{\mathrm{out}})
\;\gtrsim\;
C'\sqrt{n p_{\max}\log n},
\]
where $p_{\max}=\max\{p_{\mathrm{in}},p_{\mathrm{out}}\}$. 
% \lc{We might change $p_{\max}=p_{\mathrm{in}}$ as we assume $p_{\mathrm{in}} > p_{\mathrm{out}}$}
\end{theorem}

\begin{proof}
\label{proof_theorem}
First, we want to relate the operator-level perturbation bound from Lemma~\ref{lem:resolvent} to the per-node (entrywise) deviation of the PPR vector.
Let
\[
\Delta r := (I - \alpha P)^{-1}\mathbf{e}_{s_0} - (I - \alpha \bar P)^{-1}\mathbf{e}_{s_0},
\]
which captures the difference between the random PPR vector and its mean-field counterpart.
By Lemma~\ref{lem:resolvent}, we have the operator-level bound
\[
\|(I - \alpha P)^{-1} - (I - \alpha \bar P)^{-1}\|_2 \;\le\; \frac{2\alpha\varepsilon_P}{\delta^2},
\]
under the assumptions $\alpha\|\bar P\| \le 1-\delta$ and $\alpha\varepsilon_P \le \delta/2$. To obtain an entrywise guarantee, observe that
\[
\|\Delta r\|_\infty
= \|[(I - \alpha P)^{-1} - (I - \alpha \bar P)^{-1}]\,\mathbf{e}_{s_0}\|_\infty
\;\le\;
\|(I - \alpha P)^{-1} - (I - \alpha \bar P)^{-1}\|_{1\to\infty}.
\]
Using the inequality $\|A\|_{1\to\infty} \le \sqrt{n}\,\|A\|_2$,
we obtain the conservative bound
\[
\|\Delta r\|_\infty
\;\le\; \sqrt{n}\cdot\frac{2\alpha\varepsilon_P}{\delta^2}
\;=:\; \varepsilon_{\infty}.
\]
That is, with high probability, each node's PPR score deviates from its mean-field prediction
by at most $\varepsilon_{\infty}$ in absolute value.

\vspace{3mm}
\noindent Next, we want to prove that if $\Delta_{\mathrm{mean}} \;\ge\; C\cdot\varepsilon_{\infty},$ for a constant $C>0$ with probability at least $1-n^{-c}$, then (1) and (2) hold. 
We assume that this mean-field gap dominates the perturbation magnitude:
\[
\Delta_{\mathrm{mean}} \ge C\,\varepsilon_{\infty},
\]
for a sufficiently large absolute constant $C>0$.
Under this assumption, by concentration inequalities (e.g., Markov or Hoeffding combined with a union bound),
the fraction of nodes whose PPR values deviate by more than $\varepsilon_{\infty}$ from their mean-field predictions
can be made arbitrarily small.
Consequently, with probability at least $1 - n^{-c}$ for some constant $c>0$, we have
\[
r(v) \ge \mu_{\mathrm{in}} - \varepsilon_{\infty}
\quad \text{for most } v \in S^\star,
\qquad\text{and}\qquad
r(u) \le \mu_{\mathrm{out}} + \varepsilon_{\infty}
\quad \text{for most } u \notin S^\star.
\]
where $\mu_{\mathrm{in}}$ is the expected PPR score for a node inside the planted laundering group $S^\star$ (in-group) and $\mu_{\mathrm{out}}$ is the expected PPR score for a node outside the planted group (out-group) defined in Equation \ref{mu_in_and_out}.
To establish part (1), we require the lowest in-group score to be sufficiently separated from the highest out-group score. We need
\[
\Delta_{\mathrm{mean}} - 2\varepsilon_{\infty} \ge \frac{1}{2}\Delta_{\mathrm{mean}},
\]
which simplifies to $\Delta_{\mathrm{mean}} \ge 4\varepsilon_{\infty}$.
By choosing $C \ge 4$, the assumed condition $\Delta_{\mathrm{mean}} \ge C \varepsilon_{\infty}$ is satisfied, guaranteeing:
\begin{align*}
\text{Avg}(S^\star) - \text{Avg}(V \setminus S^\star) &\ge (\mu_{\mathrm{in}} - \varepsilon_{\infty}) - (\mu_{\mathrm{out}} + \varepsilon_{\infty}) \\
&= \Delta_{\mathrm{mean}} - 2\varepsilon_{\infty} \\
&\ge \Delta_{\mathrm{mean}} - \frac{1}{2}\Delta_{\mathrm{mean}} = \frac{1}{2}\Delta_{\mathrm{mean}}.
\end{align*}
This establishes part (1) of the theorem.

% When $C$ is large enough, these inequalities ensure that
% \[
% \mu_{\mathrm{in}} - \varepsilon_{\infty} > \mu_{\mathrm{out}} + \varepsilon_{\infty},  \Longrightarrow \mu_{\mathrm{in}} - \mu_{\mathrm{out}} > 2\varepsilon_{\infty}.
% \]
% Thus, a majority of in-group nodes obtain larger PPR values than a majority of out-group nodes.
% This establishes part (1) of the theorem: the average normalized PPR score of the planted group
% exceeds that of the rest by at least $\tfrac{1}{2}\Delta_{\mathrm{mean}}$ with $C\geq 4$.

\vspace{3mm}
\noindent Next, we argue that sorting nodes by their normalized PPR scores $r(v)/d(v)$ and performing a conductance-based sweep ~\parencite{andersen2006local} recovers a constant fraction of the true planted subset. Intuitively, since PPR mass decays smoothly away from the seed node, nodes within the planted group $S^\star$ have substantially higher normalized scores under the separation established. When the nodes are ordered by decreasing $r(v)/d(v)$, most of the top-ranked nodes belong to $S^\star$, and the boundary between $S^\star$ and $V\setminus S^\star$ forms a low-conductance cut.
By the guarantees of the local clustering framework~\parencite{andersen2006local}, sweeping over prefixes of this ordering and selecting the prefix with minimal conductance yields a subset $\widehat S$ whose conductance is within a constant factor of the optimal local set around the seed. Therefore, $\widehat S$ necessarily overlaps with $S^\star$ on at least a constant fraction, i.e.,
\[
    |\widehat S \cap S^\star| \ge \gamma s,
\]
for some constant $\gamma \in (0,1)$ depending only on the signal-to-noise separation and conductance gap. This establishes part (2) of the theorem.

\vspace{3mm}
\noindent To derive the explicit detectability threshold, we approximate the mean-field separation $\Delta_{\mathrm{mean}}$
and the perturbation scale $\varepsilon_{\infty}$.
For a two-block SBM with $s \ll n$ and $p_{\mathrm{out}}\ll 1$, expanding $(I - \alpha M)^{-1}$ shows that
\[
\Delta_{\mathrm{mean}} \asymp \frac{s(p_{\mathrm{in}} - p_{\mathrm{out}})}{n}.
\]
Substituting the estimate
$\varepsilon_P = C_1\frac{\sqrt{n p_{\max}\log n}}{n p_{\mathrm{out}}}$
from Lemma~\ref{lem:transition} into Equation \eqref{eq:epsinf}
gives
\[
\varepsilon_{\infty}
\;\asymp\;
\sqrt{n}\cdot\frac{\alpha}{\delta^2}\cdot
\frac{\sqrt{n p_{\max}\log n}}{n p_{\mathrm{out}}}
\;=\;
\frac{\alpha}{\delta^2}\cdot
\frac{\sqrt{n p_{\max}\log n}}{p_{\mathrm{out}}\sqrt{n}}.
\]
% Balancing $\Delta_{\mathrm{mean}}$ and $\varepsilon_{\infty}$ up to constants $C'$
% yields the informal detectability threshold
% \[
% s(p_{\mathrm{in}} - p_{\mathrm{out}})
% \;\gtrsim\;
% C'\sqrt{n p_{\max}\log n},
% \]
Balancing the signal $\Delta_{\mathrm{mean}}$ and $\varepsilon_{\infty}$, up to constants $C''$ dependent on $\alpha$ and $\delta$:
\[
\frac{s(p_{\mathrm{in}} - p_{\mathrm{out}})}{n}
\;\gtrsim\;
C''\frac{\sqrt{n p_{\max}\log n}}{p_{\mathrm{out}}\sqrt{n}}.
\]
Rearranging to isolate the signal terms on the left side yields the formal detectability threshold:
\[
s(p_{\mathrm{in}} - p_{\mathrm{out}})
\;\gtrsim\;
C''\frac{n}{p_{\mathrm{out}}\sqrt{n}} \sqrt{n p_{\max}\log n}
\;=\;
C'\sqrt{np_{\max}\log n}.
\]
% for a constant $C'$ depending on $\alpha$, $\delta$, and concentration constants.
The condition provided in the theorem statement,
$s(p_{\mathrm{in}}-p_{\mathrm{out}}) \;\gtrsim\; C'\sqrt{n p_{\max}\log n}$, is obtained by absorbing the factor $\frac{n}{p_{\mathrm{out}}\sqrt{n}}$ into the constant $C'$. This condition expresses that the effective signal strength of the planted community must exceed the level of random fluctuations induced by graph noise for PPR-based detection to succeed. Combining the above arguments completes the proof.


\end{proof}
